\documentclass[10pt,conference]{IEEEtran}
\usepackage{cite}
\usepackage{amsmath,amssymb}
\usepackage{booktabs}
\usepackage{array}
\usepackage{url}
\usepackage[T1]{fontenc}
\renewcommand{\arraystretch}{1.08}
\begin{document}
\title{Will It Be Possible to Cure All Cancers?\\A Design-Only Framework for Testing Broad Immune Recognition and Durable Cancer Control}
\author{Survey--Idea--ExperimentDesign--Author Research Workflow}
\maketitle
\begin{abstract}
``Cure all cancers'' is not one biomedical endpoint because cancers differ by tissue, genotype, antigen presentation, immune microenvironment, metastatic route, treatment history, and evolutionary capacity. Cancer immunotherapy has established that immune mechanisms can yield major and sometimes durable benefit in specific clinical contexts, but no discovery or current treatment proves universal cancer eradication. A reported MR1-restricted T-cell receptor (TCR) cancer-recognition strategy is an important discovery-stage hypothesis; laboratory killing of selected cancer cells does not establish tumor coverage, normal-tissue safety, persistence, resistance control, or clinical benefit. This design-only paper defines Pan-cancer Antigen and Networked Cure Readiness Evaluation (PAN-CURE), a staged framework for assessing a candidate immune-recognition strategy. PAN-CURE freezes an indication set, target/normal-tissue boundary, mechanism, delivery/product definition, resistance family, and endpoint hierarchy. It requires independently auditable target measurement, target-dependence tests, heterogeneous tumor and microenvironment panels, escape analyses, noncompensatory normal-tissue and systemic safety gates, staged clinical evaluation, relapse follow-up, and replication. The paper reports no experiments, tumor killing, response, survival, toxicity, or cure result. Its conclusion is conditional: broad immune-recognition strategies can be investigated systematically, but a credible claim must be indication-specific and must survive heterogeneity, escape, safety, and durability tests; no current evidence establishes that all cancers can be cured by one therapy.
\end{abstract}
\begin{IEEEkeywords}cancer immunotherapy, pan-cancer therapy, tumor heterogeneity, immune escape, MR1, T-cell receptor, CAR-T, resistance, clinical translation.\end{IEEEkeywords}
\section{Introduction}\label{sec:intro}
Cancer is a disease in which abnormal cells can grow uncontrollably and spread to other parts of the body~\cite{nci}. This broad definition already explains why a universal cure claim is demanding. Solid tumors and hematologic malignancies differ in location, cellular origin, genomic alterations, stromal structure, antigen presentation, immune infiltration, metastatic behavior, and treatment exposure. A treatment may be curative in a local disease, useful for a molecular subtype, inactive in another subtype, or effective initially but followed by resistant relapse. The scientific task is not to aggregate these differences into a rhetorical average; it is to state the conditions under which a specific intervention is expected to work and to reveal conditions where it fails.

Cancer biology also creates a moving target. Cancer cells can evade immune detection, reshape immune cells, alter nutrient use, recruit vasculature, accumulate chromosome alterations, invade, and metastasize~\cite{nci,hanahan2022}. An immune-recognition strategy therefore needs more than a target detected in an initial tumor panel. It must show that the recognition mechanism is available across the declared cancer states, remains available under pressure, produces killing rather than a laboratory artifact, reaches tumor sites, survives or functions in suppressive microenvironments, and does not cause unacceptable normal-tissue injury or systemic toxicity.

The prompt highlights an MR1-related TCR discovery. Crowther \emph{et al.}\ used genome-wide CRISPR screening to investigate cancer targeting through monomorphic MHC class I-related protein MR1~\cite{crowther2020}. This motivates a careful hypothesis: a metabolically altered cancer state might create an immune-recognition opportunity that is not tied to a single conventional polymorphic antigen. It does not imply that MR1 expression itself marks all cancers, that every cancer presents the same relevant ligand, that normal tissues are unaffected in humans, or that a T-cell product produces durable clinical cures. The gap between discovery, preclinical mechanism, product qualification, and human benefit is central to this report.

We introduce \emph{PAN-CURE}, Pan-cancer Antigen and Networked Cure Readiness Evaluation. PAN-CURE is not a therapy and does not rank treatments by a universal score. It defines a staged evidence architecture for a candidate broad-recognition strategy. It separates target/normal-tissue measurement, mechanistic causality, tumor heterogeneity and escape, product and systemic safety, clinical activity, durable disease control, and replication. This is a design-only paper; it reports no participant, animal, cell, tumor, or clinical result.

\section{Survey: Why a Universal Cancer-Cure Claim Is Difficult}\label{sec:survey}
\subsection{Cancer heterogeneity and endpoint definitions}
Cancer cure can mean eradication after local treatment, sustained absence of detectable disease, long-term remission, or a population-level outcome defined after sufficient follow-up. These endpoints are not interchangeable. Imaging, molecular residual disease, symptoms, survival, and recurrence each have measurement limits. A short tumor shrinkage observation cannot establish a durable cure, and a biomarker response cannot replace a patient-relevant endpoint without validation.

For individual $i$, an event-free probability through a stated horizon $T$ is
\begin{equation}
S_i(T\mid x_i,a_i)=\Pr(E_i(T)=0\mid x_i,a_i),\label{eq:eventfree}
\end{equation}
where $E_i(T)$ denotes the protocol-defined event process, $x_i$ is the declared disease state, and $a_i$ is assignment or exposure. Equation~\eqref{eq:eventfree} emphasizes that durability is conditional on disease definition, treatment context, follow-up, and censoring rules. It is not an all-cancer cure probability.

Hanahan's updated hallmarks framework describes cancer as a set of acquired capabilities with enabling features and emerging dimensions~\cite{hanahan2022}. The framework is useful for resistance hypotheses because tumor proliferation, cell death, metabolism, immune evasion, angiogenesis, invasion, and genome instability can interact. It is not a promise that a single intervention will control every capability in every tumor.

\begin{table*}[!t]\caption{Evidence boundary for a broad immune-recognition strategy.}\label{tab:boundary}\centering\footnotesize
\begin{tabular}{@{}p{.18\textwidth}p{.25\textwidth}p{.26\textwidth}p{.23\textwidth}@{}}\toprule
Evidence class & What it can establish & Main limitations & Unsupported conclusion\\\midrule
Cell-culture recognition or killing & A mechanism hypothesis in stated lines and assay conditions. & Target density, culture adaptation, missing microenvironment, donor and assay effects. & Patient efficacy, safety, or universal coverage.\\
Target/ligand atlas & Distribution evidence for specified tumors and normal samples. & Sampling, subtype coverage, dynamic presentation, spatial and treatment effects. & Stable target availability in every cancer.\\
Heterogeneous model panel & Coverage and escape map across declared models. & Model transfer, selection, limited rare states, and incomplete immune context. & Clinical durability or absence of resistance.\\
Early clinical study & Safety/activity signal for enrolled population and protocol. & Small samples, follow-up, comparator, selection, and competing therapies. & Cure of other cancers or lifelong disease eradication.\\
Replication and long follow-up & Reproducibility and durability for stated indications. & Finite horizon and population specificity. & One treatment cures all present and future cancers.\\\bottomrule\end{tabular}\end{table*}

\subsection{Immune recognition and escape}
Immune therapies can change the treatment landscape, but immune response is a system property. Checkpoint inhibition, engineered cellular products, cytokine approaches, vaccines, and other modalities differ in target, delivery, persistence, toxicity, and disease setting. Pardoll reviewed immune checkpoint blockade as a mechanism that can release anti-tumor responses, while also underscoring the logic of tumor-immune interaction~\cite{pardoll2012}. A therapy's activity may depend on antigen presentation, T-cell trafficking, tumor burden, suppressive myeloid and stromal states, prior therapy, and host factors.

Target availability is usefully treated as a distribution, not a yes/no label. Let $Z_{cg}$ represent a target or recognition-relevant measurement for cancer state $c$ in context $g$, and let $N_g$ represent a matched normal-tissue measurement. A declared coverage-and-safety margin is
\begin{equation}
M_{cg}=Z_{cg}-\max_{g'\in\mathcal{N}}N_{g'},\label{eq:margin}
\end{equation}
where $\mathcal{N}$ is the predeclared normal-tissue set. Equation~\eqref{eq:margin} is a measurement rule, not a safety certificate: expression or presentation values require assay calibration, spatial context, and functional testing. A positive average margin can still hide a normal cell subset or tumor clone that determines toxicity or escape.

MR1 is an especially instructive example. A monomorphic molecule could, in principle, simplify one dimension of recognition compared with highly polymorphic restrictions. But an MR1-restricted response still requires mechanistic evidence about the relevant ligand, tumor-state dependence, normal-tissue boundary, and selection under treatment pressure. PAN-CURE treats MR1/TCR strategies as candidates for S0--S2 evaluation, not as evidence of a ``one-stop'' treatment.

\subsection{Resistance, metastasis, and the tumor microenvironment}
Therapy exerts selection. A tumor can lose or reduce a target, alter processing/presentation, change lineage or metabolic state, exclude immune cells, suppress effector function, or activate compensatory survival pathways. Resistant clones may pre-exist or emerge during treatment. Metastatic sites can have distinct microenvironments and antigen states. A strategy that succeeds in one lesion or model can fail elsewhere.

For a declared resistance family $\mathcal{R}$, define the conditional escape probability
\begin{align}
p_{\mathrm{esc}}(T\mid\mathcal{R},x,a)&=\Pr\left(\exists r\in\mathcal{R}:r\ \mathrm{persists}\right.\nonumber\\
&\left.\quad\mathrm{or\ expands\ by}\ T\mid x,a\right).\label{eq:escape}
\end{align}
Equation~\eqref{eq:escape} is not estimated in this paper. It states the parameter a future program must bound with molecular, functional, and longitudinal evidence before using broad durability language.

\section{Idea: PAN-CURE}\label{sec:idea}
PAN-CURE converts the question into a bounded development claim: for an immune modality, indication set, target mechanism, product definition, and time horizon, does the strategy show reproducible selective activity, resistance-aware coverage, acceptable safety, and durable patient-relevant benefit? The null hypothesis $H_0$ attributes apparent breadth to selected models, target overexpression, off-target assay effects, reporting bias, or inadequate follow-up. Under $H_1$, the intervention passes the preregistered gates for a declared cancer set. Under $H_2$, activity is restricted to a subset or context. Under $H_3$, toxicity, escape, or microenvironmental failure blocks the claim.

The core idea is noncompensation. Strong in-vitro killing cannot offset credible normal-tissue injury. Response in one subtype cannot offset antigen-negative or resistant outgrowth. A favorable short-term response cannot offset absent durability. Each claim has a boundary: tumor types and molecular states, disease stage, prior therapy, immune context, target assay, product/delivery, comparator, outcome, and follow-up.

For a multiparameter evidence vector $V=(C,M,R,S,D)$ representing coverage, mechanism, resistance, safety, and durability, PAN-CURE reports component values rather than a hidden aggregate. A logical readiness gate is
\begin{equation}
G=\mathbb{I}_{C}\mathbb{I}_{M}\mathbb{I}_{R}\mathbb{I}_{S}\mathbb{I}_{D},\label{eq:readiness}
\end{equation}
where every indicator must meet a preregistered criterion. Equation~\eqref{eq:readiness} ensures that $G=0$ whenever a coverage, mechanism, resistance, safety, or durability condition fails.

\begin{table*}[!t]\caption{PAN-CURE endpoint hierarchy.}\label{tab:endpoints}\centering\footnotesize
\begin{tabular}{@{}p{.18\textwidth}p{.27\textwidth}p{.25\textwidth}p{.22\textwidth}@{}}\toprule
Tier & Required measurement & Inferential role & Prohibited shortcut\\\midrule
Target boundary & Tumor subtype, normal tissue, spatial and longitudinal target/ligand measurements. & Defines feasible recognition boundary. & Calling a molecule pan-cancer from limited samples.\\
Mechanism & Target dependence, receptor specificity, ablation/restoration, functional immune readouts. & Tests causal recognition mechanism. & Inferring mechanism from correlation or killing alone.\\
Heterogeneity/escape & Diverse models, low-target states, microenvironment challenges, resistant-outgrowth analyses. & Maps coverage limits and escape. & Replacing an escape map with mean response.\\
Safety & Normal-tissue comparators, cytokine/systemic risk plan, product identity and persistence controls. & Determines whether efficacy interpretation is eligible. & Offsetting safety concern by response.\\
Clinical durability & Protocol-defined activity, toxicity, relapse, survival/event follow-up, replication. & Tests patient-relevant benefit for stated indication. & Generalizing a short response to universal cure.\\\bottomrule\end{tabular}\end{table*}

\section{ExperimentDesign: Staged Broad-Recognition Evaluation}\label{sec:design}
\subsection{Scenario ladder and controls}
S0 creates an audited target/normal-tissue and tumor-state atlas. S1 tests recognition and target dependence with matched target-negative, irrelevant-receptor, and normal comparators. S2 applies a heterogeneous panel incorporating subtypes, target density, antigen-presentation states, treatment history where available, stromal/myeloid context, and predefined resistance challenges. S3 is a translational safety qualification whose systems depend on the candidate modality. S4 is a staged clinical program with ethics, regulator, and safety oversight. No lower rung licenses a higher-rung conclusion.

The primary source/provenance ledger records specimen state, assay, batch, target/ligand measurement, receptor/product version, culture or model condition, treatment exposure, and analysis lock. Target ablation and restoration, where appropriate, test dependence. Multiple donors and technical controls test whether a response is driven by one effector preparation. Blinded labels and held-out models prevent a target hypothesis from being adjusted after response inspection.

For model $m$ and time $t$, a bounded tumor-to-normal selectivity contrast can be written
\begin{equation}
\Delta_{m}(t)=\mathbb{E}[Y_{\mathrm{tumor},m}(t)\mid A=1]-\mathbb{E}[Y_{\mathrm{normal},m}(t)\mid A=1],\label{eq:selectivity}
\end{equation}
where $Y$ is a preregistered functional readout and $A$ denotes intervention exposure. Equation~\eqref{eq:selectivity} requires exact assay definition and does not itself establish safe human selectivity.

\begin{table*}[!t]\caption{PAN-CURE scenario ladder and permissible conclusions.}\label{tab:ladder}\centering\footnotesize
\begin{tabular}{@{}p{.15\textwidth}p{.29\textwidth}p{.28\textwidth}p{.20\textwidth}@{}}\toprule
Scenario & Required evidence & Failure disqualifier & Permissible conclusion\\\midrule
S0 atlas & Sample provenance, tumor/normal coverage, subtype and spatial context, assay controls. & Normal boundary or clinically relevant tumor states are omitted. & Measurement boundary is defined.\\
S1 mechanism & Target dependence, matched controls, normal comparators, receptor/product identity, replicate donors. & Killing persists without target mechanism or controls fail. & Bounded recognition mechanism is supported.\\
S2 heterogeneous models & Predeclared subtype/low-target/microenvironment/resistance panel and null reporting. & Coverage inferred from selected sensitive models. & Model-specific coverage and escape map is supported.\\
S3 safety qualification & Candidate-specific toxicology, persistence, biodistribution, cytokine and off-tumor plans. & Credible safety issue or incomplete product characterization. & Preclinical safety package under stated limits.\\
S4 clinical program & Endpoint, safety adjudication, relapse, censoring, molecular escape, and follow-up plan. & Endpoint switching, inadequate follow-up, or concealed toxicity. & Indication-specific clinical evidence.\\\bottomrule\end{tabular}\end{table*}

\subsection{Heterogeneity, escape, and safety}
Heterogeneity must be built into the test set rather than treated as a post hoc explanation. The panel is frozen before analysis and includes target-low or target-variable states, distinct molecular subtypes, relevant metastatic contexts where models allow, and suppressive microenvironment conditions. Resistance experiments test target loss, pathway alteration, antigen-presentation changes, and phenotype/metabolic shifts appropriate to the recognition hypothesis. A negative model is not discarded merely because it contradicts a broad narrative.

Safety requires an independent boundary. Normal tissue can share molecules, stress programs, or presentation context with tumors. Engineered immune products can also generate systemic inflammation, prolonged persistence, or other modality-specific risks. The study plan includes product identity, contamination control, functional normal comparators, dose/exposure logic for its model, cytokine and systemic monitoring plans, biodistribution and persistence assessment where applicable, and explicit pause/stop rules for later-stage work. None of these assays has been performed here.

Clinical durability is modeled separately from response. With event indicator $E_i$ and retained follow-up indicator $R_i$, an inverse-probability sensitivity weight is
\begin{equation}W_i(t)=\prod_{s\leq t}\frac{q_s(1\mid A_i)}{q_s(1\mid A_i,L_i(s))},\label{eq:retention}\end{equation}
where $L_i(s)$ denotes preregistered covariate, disease, and treatment history. Equation~\eqref{eq:retention} is only a sensitivity model; it does not solve informative censoring or convert short follow-up into cure evidence.

\section{Microenvironment, Delivery, and Clinical Estimands}\label{sec:clinical}
\subsection{A recognition mechanism is not a complete treatment system}
An immune receptor acts within a treatment system. The same recognition event can have different consequences depending on trafficking, immune-cell fitness, tumor burden, stromal exclusion, suppressive myeloid programs, hypoxia, nutrient competition, cytokine signals, prior treatment, and anatomical site. A product can fail to reach a metastasis, reach it but be inactivated, or generate an inflammatory response that narrows the therapeutic window. These are not secondary engineering details; they determine whether a receptor-level discovery can become an indication-level intervention.

PAN-CURE therefore keeps the model context explicit. A model panel includes not only target-high cell lines but target-low and variable states, relevant molecular subtypes, altered antigen-presentation states, and microenvironmental conditions that challenge effector function. When a model is simplified, the conclusion is labeled as such. A successful recognition assay in an isolated culture can support a mechanism hypothesis, but it cannot answer trafficking, persistence, host toxicity, or relapse. A negative result in a suppressive model is not discarded as noise; it may identify a necessary combination strategy or an indication boundary.

The decision-relevant question is whether an additional experiment changes a gate. For example, a broader tumor atlas matters if it reveals that the candidate mechanism is absent in clinically relevant states. A matched normal-tissue assay matters if it changes the plausible on-target/off-tumor boundary. A resistance experiment matters if it shows that target loss or pathway adaptation can defeat the approach. Repeating only sensitive models cannot close any of these dependencies.

\begin{table*}[!t]\caption{Context variables that can turn an immune-recognition finding into success, resistance, or toxicity.}\label{tab:context}\centering\footnotesize
\begin{tabular}{@{}p{.18\textwidth}p{.26\textwidth}p{.27\textwidth}p{.21\textwidth}@{}}\toprule
Context variable & Required evidence & Failure mode if omitted & Bounded inference enabled\\\midrule
Tumor-state and target dynamics & Longitudinal and subtype-resolved target/ligand, presentation, and pathway measurements. & A static target atlas misses treatment-induced or resistant states. & Recognition boundary for stated samples and times.\\
Immune trafficking and fitness & Location, infiltration, persistence, activation/exhaustion, and local suppressive-context measurements. & Receptor activity in vitro is mistaken for activity at tumor sites. & Mechanistic feasibility in the stated model.\\
Stroma and myeloid environment & Matrix, cytokine, macrophage/myeloid, and exclusion conditions matched to the indication where feasible. & An accessible culture target is assumed to be accessible in disease. & Microenvironment sensitivity or combination need.\\
Normal tissue and systemic exposure & Matched normal comparators, product distribution/persistence and modality-specific risk plan. & Mean tumor selectivity hides a rare but clinically important toxic tissue. & Restricted safety hypothesis, not human safety proof.\\
Evolution under treatment & Target-loss, presentation-loss, phenotype, metabolic, and clonal-outgrowth stress tests. & Initial activity is misreported as durable control. & Preclinical escape map and follow-up priorities.\\\bottomrule\end{tabular}\end{table*}

\subsection{Clinical endpoints and the meaning of cure}
Clinical translation needs an estimand before enrollment. The target population, disease stage, molecular subgroup, prior treatment, product exposure, comparator, rescue therapy, adverse-event window, and follow-up must be fixed. Safety, objective activity, molecular residual disease, progression, relapse, survival, patient-reported outcomes, and quality of life have distinct roles. A response endpoint can be important for early development but does not prove durable eradication. A relapse endpoint requires complete surveillance and molecular characterization of recurrence. Survival requires careful handling of subsequent therapy and censoring.

Let $Y_i(t)$ be a protocol-defined longitudinal tumor or molecular measurement and let $H_i(t)$ denote safety history. A prospective analysis can separate short-term activity from durable benefit using
\begin{equation}
\Delta_Y(t)=\mathbb{E}[Y_i(t)\mid A_i=1]-\mathbb{E}[Y_i(t)\mid A_i=0],\label{eq:activity}
\end{equation}
while reporting $H_i(t)$ and the event process in Equation~\eqref{eq:eventfree} separately. Equation~\eqref{eq:activity} is not a cure criterion. A favorable $\Delta_Y(t)$ can coexist with later relapse, toxicity, or unequal follow-up.

The clinical program should also preserve molecular escape analysis. A recurrence is not merely a failed outcome; it can test whether the target or presentation pathway changed, whether a different clone expanded, or whether an immune microenvironmental barrier dominated. These analyses are hypothesis-driven and require consent, assay validation, and predeclared interpretation rules. They should not be used to relabel a nonresponse as success after the fact.

\begin{table*}[!t]\caption{Clinical evidence ladder for an indication-specific immune strategy.}\label{tab:clinical-ladder}\centering\footnotesize
\begin{tabular}{@{}p{.17\textwidth}p{.27\textwidth}p{.27\textwidth}p{.21\textwidth}@{}}\toprule
Stage & Primary question & Mandatory safety/durability information & Conclusion permitted\\\midrule
Discovery and preclinical mechanism & Is the recognition mechanism target dependent in the declared model? & Normal comparators, assay controls, model limitations, and escape hypotheses. & Mechanistic candidate, not a therapy.\\
Translational qualification & Is a defined product/exposure interpretable and sufficiently characterized for the next study? & Product identity, distribution/persistence, toxicology plan, stop rules, and manufacturing controls. & Candidate for supervised clinical evaluation.\\
Early clinical study & What safety and activity signals occur in the enrolled indication and protocol? & Adverse events, dose/exposure provenance, complete response/nonresponse reporting, and follow-up. & Population-specific early clinical evidence.\\
Comparative/durability study & Does the strategy improve prespecified patient-relevant outcomes versus the comparator? & Relapse, subsequent therapy, survival/event follow-up, patient outcomes, molecular escape, and replication. & Bounded durable-benefit claim for stated indication.\\
Cross-indication extension & Does the mechanism and benefit reproduce under a separately specified cancer and patient boundary? & Independent target/normal, microenvironment, resistance, and safety evaluation. & Evidence for an additional indication, not universal cure.\\\bottomrule\end{tabular}\end{table*}

\subsection{Combinations and interaction claims}
Combination treatments may be scientifically justified when a receptor-level strategy is constrained by a separable barrier, such as trafficking, antigen presentation, suppression, or a resistance route. But a combination raises new attribution and safety questions. If exposures $A$ and $B$ are tested, the protocol distinguishes the effect of the full strategy from a component interaction; it does not infer synergy merely because a combination seems stronger than one selected control. A transparent interaction contrast is
\begin{equation}
I_{AB}=\Delta_{AB}-\Delta_A-\Delta_B,\label{eq:interaction}
\end{equation}
where all contrasts use the same locked model, endpoint, and follow-up. Equation~\eqref{eq:interaction} is descriptive unless the experimental design and scale justify a causal interaction interpretation. It also does not compensate for combined toxicity.

\section{Mechanism Confirmation and Research Priorities}\label{sec:priorities}
An MR1-restricted TCR or another broad-recognition candidate should be developed as a falsifiable mechanism program. First, the molecular recognition condition must be narrowed: which tumor-associated state, ligand or presentation configuration is necessary; whether target ablation removes recognition; whether restoration restores it; and whether the same dependence holds in independent tumor contexts. Second, the normal-tissue boundary must be mapped using the same assays and exposure conditions as the tumor comparison. Third, the hypothesis must face adversarial tumor states, including low-recognition, altered-presentation, stress, metabolic, and resistant states. A result that remains strong only after those states are excluded is a subset result, not a pan-cancer result.

The order of these tests matters. A large clinical expansion cannot repair an unresolved target mechanism or a plausible normal-tissue hazard. Conversely, a mechanistic failure does not invalidate the general field of immunotherapy; it identifies the specific dependency that blocked the candidate. This is why PAN-CURE reports a dependency ledger rather than an average efficacy score. The ledger helps distinguish a target-distribution problem, a receptor/product problem, a microenvironment problem, an evolution problem, and a safety problem.

\begin{table*}[!t]\caption{Research priorities for a broad immune-recognition hypothesis.}\label{tab:priorities}\centering\footnotesize
\begin{tabular}{@{}p{.19\textwidth}p{.26\textwidth}p{.27\textwidth}p{.20\textwidth}@{}}\toprule
Open question & Discriminating experiment or data & Decision changed by the result & Result does not establish\\\midrule
What is recognized? & Target/ligand dependence, receptor mutation, ablation/restoration, and independent tumor states. & Whether the candidate mechanism is causal and how to define target-positive disease. & Human safety or clinical response.\\
How broad is tumor coverage? & Frozen subtype, stage, metastatic-site, target-low, and treatment-history atlas. & Whether the indication should narrow, expand, or stop. & Durable benefit across patients.\\
What normal tissue is vulnerable? & Matched normal-tissue and systemic-context functional tests with product/exposure controls. & Whether a safety gate blocks further translation. & Absence of rare or long-term toxicity.\\
How does escape occur? & Predeclared pressure, target/presentation-loss, phenotype, metabolic, and clonal-outgrowth tests. & Whether resistance mitigation or combinations are necessary. & Prevention of relapse in people.\\
Which patients benefit? & Prospectively defined molecular and immune-state stratification with complete nonresponder reporting. & Whether a clinical population can be specified. & Efficacy outside the validated population.\\\bottomrule\end{tabular}\end{table*}

Patient stratification is not a failure of the cure ambition. It is an empirical response to biological variation. The most ethical and informative clinical question is often not ``does the product work?'' but ``for which molecular, immune, disease-stage, and prior-treatment boundary does the benefit-risk profile support use?'' A biomarker hypothesis used for enrichment must be locked and evaluated for both sensitivity and exclusion risk. Excluding a patient group based on an unvalidated marker can create a false safety or efficacy narrative; including all groups without measurement can dilute an interpretable signal.

The same care applies to relapse. Molecular changes at recurrence can indicate an escape route, but recurrent sampling has selection and feasibility limitations. PAN-CURE records which samples were available, which were not, and whether recurrence analyses were predefined. A resistant clone discovered after treatment is evidence about a selected outcome, not proof that the same route would dominate in all cancers. The framework turns such observations into next-stage hypotheses rather than post hoc claims of mechanism.

\section{Communication and Translation Boundaries}\label{sec:communication}
The phrase ``cure all cancers'' can create unjustified expectations, especially for people facing serious disease. Research communication should name the experimental system, the candidate target, the cancer boundary, the response and safety endpoint, and the remaining failure modes. It should distinguish an immune-cell laboratory observation from a manufactured cell product, a preclinical model from a patient, and an early response from a durable patient-relevant outcome. This is not merely cautious wording: these differences determine the design, oversight, and evidence required for the next step.

PAN-CURE requires a claim ledger for every release. The ledger states the observation, mechanism inference, preclinical interpretation, or clinical result; all omitted cancer states; all normal-tissue assumptions; resistance tests; toxicity observations; and follow-up horizon. It also records null, resistant, and toxic results. This practice reduces the chance that a striking example of tumor killing becomes an implied promise of universal therapeutic success.

\section{Interface Failure Modes}\label{sec:interfaces}
Broad cancer-control programs fail most often at interfaces rather than within a single assay. A receptor can recognize a target but fail to persist; a product can persist but fail to traffic; immune cells can enter a lesion but face local suppression; a tumor can be initially sensitive but evolve an escape state; a trial can show activity but lack the follow-up needed to distinguish remission from durable control. PAN-CURE assigns each interface an owner, measurement, contingency, and claim boundary. This makes negative results informative and prevents a strong component result from being represented as end-to-end efficacy.

\begin{table*}[!t]\caption{Interface failure matrix for a broad immune-recognition program.}\label{tab:interfaces}\centering\footnotesize
\begin{tabular}{@{}p{.18\textwidth}p{.25\textwidth}p{.28\textwidth}p{.21\textwidth}@{}}\toprule
Interface & Example failure & Evidence required before stronger claim & Consequence for interpretation\\\midrule
Recognition to killing & Binding or target measurement occurs without functional tumor-cell elimination. & Target-dependence, functional assays, matched controls, and time-resolved mechanism readouts. & Recognition observation only; no anti-tumor activity claim.\\
Product to tumor site & Product is active ex vivo but cannot traffic, persist, or remain functional at disease sites. & Distribution, persistence, site-context, and immune-fitness measurements appropriate to the modality. & Model mechanism does not establish in-disease feasibility.\\
Killing to durability & Initial reduction is followed by resistant or antigen-negative outgrowth. & Longitudinal molecular and event follow-up with prespecified escape analyses. & Initial activity is not durable control or cure.\\
Tumor selectivity to patient safety & Tumor-biased signal coexists with normal-tissue or systemic inflammatory injury. & Matched normal contexts, candidate-specific safety package, adjudication, and stop rules. & Benefit interpretation is blocked by safety gate.\\
Single indication to broad claim & Benefit in one cancer type is generalized despite distinct antigen, stroma, stage, or host context. & Independent target, heterogeneity, safety, and outcome evidence for each extension. & Only the tested indication claim is supported.\\\bottomrule\end{tabular}\end{table*}

This matrix is also a compact prioritization tool. When the main uncertainty is target distribution, expanding clinical enrollment is less informative than resolving the tumor-normal and escape boundary. When the main uncertainty is safety, increasingly complex efficacy models do not remove the need for a candidate-specific safety program. When the main uncertainty is durability, endpoint and follow-up design must change before language about cure changes. The framework thereby links scientific progress to the particular dependency that remains open.

\section{Systematics, Decision Gates, and Reproducibility}\label{sec:systems}
PAN-CURE uses calibration and fault tests before a broad claim. Synthetic target-positive/negative mixtures test atlas classifiers. Target-ablation labels test whether the mechanism pipeline detects loss of dependence. Batch shifts, donor effects, low-target states, and simulated resistant clones test whether an analysis reports coverage only when it is present. Null assignments test false positives. A pipeline that calls broad activity in null or target-independent controls cannot support a broad recognition claim.

Multiplicity is predeclared. For $m$ ordered confirmatory hypotheses, a family-wise step-down criterion may compare
\begin{equation}p_{(j)}\leq\frac{\alpha}{m-j+1},\label{eq:holm}\end{equation}
at step $j$. Equation~\eqref{eq:holm} makes the hypothesis family visible but does not correct selected models, flexible target definitions, or safety undermeasurement.

\begin{table*}[!t]\caption{Decision gates and value-of-information priorities.}\label{tab:gates}\centering\footnotesize
\begin{tabular}{@{}p{.18\textwidth}p{.27\textwidth}p{.27\textwidth}p{.20\textwidth}@{}}\toprule
Open dependency & Discriminating evidence & Decision changed & Claim still not proved\\\midrule
Target breadth/normal boundary & Independent tumor-normal atlas with subtype, spatial, and longitudinal context. & Whether a pan-cancer hypothesis is even testable. & Clinical selectivity or cure.\\
Recognition mechanism & Ablation/restoration, receptor controls, relevant-ligand and functional assays. & Whether to pursue the candidate mechanism. & Coverage under immune pressure.\\
Escape and microenvironment & Predeclared resistant/low-target and suppressive-context panels. & Whether combinations or narrower indications are required. & Patient durability.\\
Safety & Normal comparators and candidate-specific systemic/product risk studies. & Whether translation language is blocked. & Absence of all human risk.\\
Durability/replication & Longitudinal clinical events, relapse/escape analysis, independent confirmation. & Whether a bounded durable-benefit claim is supported. & Cure of all cancers.\\\bottomrule\end{tabular}\end{table*}

The reproducibility package releases protocol versions, model/source provenance, target and normal definitions, product/receptor version, assay and batch metadata, complete inclusion/exclusion rules, code, calibration controls, null tests, resistant models, negative results, safety observations, and an explicit distinction between observation, mechanism inference, preclinical evidence, and clinical extrapolation. In this paper no empirical package exists because no experiment was run.

\section{Threats to Validity and Claim Calibration}\label{sec:validity}
The principal threat is selection. A target can look common when samples are chosen after a promising screen, when target-high cell lines stand in for disease, or when rare normal cells and low-target tumor states are missing. A high mean signal is not a coverage estimate unless the sampling frame, assay detection limits, subtype representation, and exclusion rules are fixed before the result is examined. PAN-CURE therefore requires every apparently favorable panel to be accompanied by the denominator: which tumor states, normal contexts, donors, time points, and failed measurements were eligible but absent. This makes an unknown boundary visible rather than converting it into an implied universal one.

Transportability is a second threat. Tumor cells in culture, organoids, xenografts, and even selected early clinical populations differ in vasculature, immune composition, treatment exposure, and evolutionary history. A negative model may be uninformative when it lacks the relevant ligand or delivery condition; a positive model may be uninformative when it excludes suppressive or resistant states. The remedy is not a generic ``more models'' request. It is a prespecified challenge set whose components map to a stated failure route: antigen or presentation loss, metabolic adaptation, poor trafficking, immune suppression, product persistence, or normal-tissue exposure. Evidence can then reduce uncertainty about one route without being misdescribed as evidence against all routes.

Clinical interpretation adds further threats. Informative censoring, post-protocol therapy, unequal imaging schedules, response assessment changes, and incomplete relapse characterization can make a short observation window appear durable. A time-to-event analysis must state the event definition, time origin, censoring rule, retained follow-up, subsequent interventions, and molecular sampling plan. Safety interpretation likewise requires denominators and exposure timing, because a rare serious normal-tissue injury cannot be neutralized by a favorable average response. These conditions are not administrative details: they determine whether an activity signal can support a bounded patient-benefit statement.

Finally, replication should challenge the most attractive claim rather than merely repeat the easiest assay. Independent teams should test locked target and normal definitions, blinded or held-out samples where feasible, a different model source, and the predeclared escape and safety tests. Claim language is calibrated to the strongest replicated gate passed: discovery supports a mechanism hypothesis; preclinical evidence supports a bounded candidate rationale; clinical follow-up can support an indication-specific benefit statement. None of these labels is interchangeable with a universal cure conclusion.

\section{Discussion and Conclusion}\label{sec:discussion}
It may become possible to cure more cancers, and some cancers are already curable in particular settings, but no evidence supports a claim that one therapy will cure all cancers. The reasons are empirical rather than merely semantic: cancers are heterogeneous and evolve; immune recognition depends on target state and context; tumors can escape; microenvironments can suppress therapy; normal tissues and systemic responses constrain safety; and durable benefit requires long follow-up. A universal claim must solve all of these simultaneously.

PAN-CURE offers a constructive route. It does not dismiss broad immune-recognition discoveries. It asks them to pass the tests required by their breadth: audited tumor and normal boundaries, causal mechanism, heterogeneous coverage, resistance, safety, clinical durability, and replication. The MR1/TCR discovery is appropriately valuable as a hypothesis that can be tested through this ladder. It is not evidence of a one-stop clinical treatment.

This four-stage design-only workflow concludes that broad cancer control should be communicated as a progressively constrained, indication-specific claim. A future program may show a defined recognition mechanism, then a coverage map, then preclinical qualification, then clinical activity and durable benefit for specific patients and cancers. Each successful gate matters; no finite set of early results can justify language that all cancers have been cured.

\appendices
\section{Minimal PAN-CURE Run Manifest}
Each future run records indication and disease-state boundary; target/ligand and normal-tissue definition; sample provenance; receptor/product version; delivery and exposure definition; target-dependence controls; model heterogeneity; microenvironment and resistance family; assay batches; safety plan; clinical endpoint, censoring, toxicity, relapse, and follow-up rules; code; failed and null results; and a statement that the run cannot certify a universal cancer cure.
\section{Protocol Checklist}
Before analysis, freeze target and normal definitions, model panel, resistance challenges, safety gates, endpoints, and follow-up. During a run, preserve target-negative, irrelevant-receptor, normal-tissue, donor, batch, and null controls; report all negative, resistant, and toxic findings. Before communication, label results as measurement, mechanism, preclinical, or clinical evidence and do not convert a finite indication-specific result into an all-cancer cure claim.
\newpage
\begin{thebibliography}{10}
\bibitem{nci}National Cancer Institute, ``What Is Cancer?'' accessed Sep. 3, 2026. [Online]. Available: \url{https://www.cancer.gov/about-cancer/understanding/what-is-cancer}
\bibitem{hanahan2022}D. Hanahan, ``Hallmarks of Cancer: New Dimensions,'' \emph{Cancer Discovery}, vol. 12, no. 1, pp. 31--46, 2022, doi: 10.1158/2159-8290.CD-21-1059.
\bibitem{crowther2020}M. D. Crowther \emph{et al.}, ``Genome-wide CRISPR-Cas9 screening reveals ubiquitous T cell cancer targeting via the monomorphic MHC class I-related protein MR1,'' \emph{Nature Immunology}, vol. 21, pp. 178--185, 2020, doi: 10.1038/s41590-019-0578-8.
\bibitem{pardoll2012}D. M. Pardoll, ``The blockade of immune checkpoints in cancer immunotherapy,'' \emph{Nature Reviews Cancer}, vol. 12, pp. 252--264, 2012, doi: 10.1038/nrc3239.
\bibitem{chen2017}D. S. Chen and I. Mellman, ``Elements of cancer immunity and the cancer-immune set point,'' \emph{Nature}, vol. 541, pp. 321--330, 2017, doi: 10.1038/nature21349.
\bibitem{ribas2015}A. Ribas and J. D. Wolchok, ``Cancer immunotherapy using checkpoint blockade,'' \emph{Science}, vol. 359, no. 6382, pp. 1350--1355, 2018, doi: 10.1126/science.aar4060.
\bibitem{june2018}C. H. June, M. Sadelain, and P. D. Greenberg, ``CAR T cell therapy,'' \emph{New England Journal of Medicine}, vol. 379, pp. 64--73, 2018, doi: 10.1056/NEJMra1706169.
\bibitem{mardis2019}E. R. Mardis, ``Neoantigens and genome instability: Impact on immunogenomic phenotypes and immunotherapy response,'' \emph{Genome Medicine}, vol. 11, Art. 71, 2019, doi: 10.1186/s13073-019-0699-2.
\bibitem{samur2019}M. K. Samur \emph{et al.}, ``Biallelic loss of BCMA as a resistance mechanism to CAR T cell therapy,'' \emph{Nature Communications}, vol. 12, Art. 868, 2021, doi: 10.1038/s41467-021-21178-7.
\end{thebibliography}
\end{document}
